\section*{第 \ref{ch:b1:acetals-protection} 章——羰基：亲核加成、缩醛与保护}

\begin{solution}{exo:b1:acetals-protection:1}
\ce{CH3CH(OH)OCH3}：\omterm{def:b1:acetals-protection:hemiacetal}{半缩醛}。\ce{CH3CH(OCH3)2}：\omterm{def:b1:acetals-protection:hemiacetal}{缩醛}。
\ce{CH3CH(OH)2}：水合物。\ce{CH3OCH3}：醚。\ce{(CH3)2C(OCH2CH3)2}：
酮的\omterm{def:b1:acetals-protection:hemiacetal}{缩醛}。
\end{solution}

\begin{solution}{exo:b1:acetals-protection:2}
\ce{CH3CH2CHO + 2CH3OH <=> CH3CH2CH(OCH3)2 + H2O}；丙酮与
乙烷-1,2-二醇反应生成 2,2-二甲基-1,3-二氧戊环和水，
\ce{(CH3)2CO + HOCH2CH2OH <=> C5H10O2 + H2O}.
\end{solution}

\begin{solution}{exo:b1:acetals-protection:3}
氢氧化钠、甲基溴化镁和硼氢化钠都不改变它；
稀硫酸水溶液可使其水解。
\end{solution}

\begin{solution}{exo:b1:acetals-protection:4}
4 号碳上的羟基加到醛基碳 1 上，形成一个五元环，
环由四个碳和一个氧组成，
与环氧相邻的碳上带有一个羟基。
\end{solution}

\begin{solution}{exo:b1:acetals-protection:5}
C=O 质子化；甲醇加到碳上；失去 \ce{H+}，得\omterm{def:b1:acetals-protection:hemiacetal}{半缩醛}；
其羟基质子化；失去水，生成 \ce{CH3CH=O+CH3}；第二个甲醇
加成；失去 \ce{H+}，得 \ce{CH3CH(OCH3)2}。被取走的两个质子又被放回，
故酸是\omterm{def:b1:elementary-steps:catalyst}{催化剂}。
\end{solution}

\begin{solution}{exo:b1:acetals-protection:6}
一个 OCH3 质子化；失去甲醇，生成 \ce{(CH3)2C=O+CH3}；水加成；
失去 \ce{H+}，得\omterm{def:b1:acetals-protection:hemiacetal}{半缩醛}；质子化、失去甲醇、再失去 \ce{H+}：
得丙酮。水过量使水解中 $Q$ 保持小于 $K$，
反应进行到底。
\end{solution}

\begin{solution}{exo:b1:acetals-protection:7}
每个\omterm{def:b1:acetals-protection:hemiacetal}{缩醛}需一份水：$0.250 \times 18.0 = \qty{4.5}{g}$，约
\qty{4.5}{mL}。一小时后，$3.2/18.0 = \qty{0.178}{mol}$：\qty{71}{\percent}
转化率。
\end{solution}

\begin{solution}{exo:b1:acetals-protection:8}
使用二元醇时，一个分子取代两个分子：反应不会使分子数减少
（一分子醛与一分子二醇生成一分子\omterm{def:b1:acetals-protection:hemiacetal}{缩醛}和一分子
水）；第二个羟基靠近反应碳，
易于成环。二者都使环状\omterm{def:b1:acetals-protection:hemiacetal}{缩醛}更有利。
\end{solution}

\begin{solution}{exo:b1:acetals-protection:9}
酸会破坏\omterm{def:b1:organomagnesium:organometallic}{格氏试剂}（质子转移）；痕量
的水和二醇也必须除去。
\end{solution}

\begin{solution}{exo:b1:acetals-protection:10}
1. 将 4-氯丁-2-酮的酮羰基以环状\omterm{def:b1:acetals-protection:hemiacetal}{缩醛}形式\omterm{def:b1:acetals-protection:protecting-group}{保护}
（乙烷-1,2-二醇、酸\omterm{def:b1:elementary-steps:catalyst}{催化剂}、迪安–斯塔克装置）。2. 在无水乙醚中与镁作用：
由\omterm{def:b1:acetals-protection:protecting-group}{保护}后的氯代物制得\omterm{def:b1:organomagnesium:organometallic}{格氏试剂}，不再被
自身羰基破坏。3. 加入丙酮，用氯化铵水溶液后处理：
得叔醇。4. 稀酸水溶液\omterm{def:b1:acetals-protection:protecting-group}{脱保护}，得
5-羟基-5-甲基己-2-酮。
\end{solution}

\begin{solution}{exo:b1:acetals-protection:11}
$Q = [\text{缩醛}][\ce{H2O}]/([\text{醛}][\ce{CH3OH}]^2)$：甲醇大量过量，
分母中呈平方，使 $Q < K$ 保持到醛
所剩无几为止。除去水作用于分子；用甲醇作溶剂较简单，
但需要易于分离的产物。
\end{solution}

\begin{solution}{exo:b1:acetals-protection:12}
只有\omterm{def:b1:acetals-protection:hemiacetal}{半缩醛}的羟基能以水的形式离去，生成一个
由相邻环氧稳定的阳离子（\ce{C=O+}）；其他羟基
须从普通碳上离去，生成不稳定的阳离子。甲醇加到
稳定阳离子上，得甲基\omterm{def:b1:acetals-protection:hemiacetal}{缩醛}。
\end{solution}

\begin{solution}{pb:b1:acetals-protection:1}
\textbf{1.} \ce{BrMgCH2CH2CH2CHO}.
\textbf{2.} 其类\omterm{def:b1:electronic-effects:intermediates}{碳负离子}样的碳会进攻另一分子
（或其自身）的醛基：\ce{RMgBr} 加到 $\mathrm{R'CHO}$ 上，生成\omterm{def:b1:alcohol-activation:alkoxide}{醇盐}。
\textbf{3.} 醛基须防\omterm{def:b1:organomagnesium:organometallic}{格氏试剂}进攻。
\textbf{4.} \omterm{def:b1:acetals-protection:hemiacetal}{缩醛}对\omterm{def:b1:organomagnesium:organometallic}{格氏试剂}和碱稳定，
可被酸的水溶液除去。
\textbf{5.} 它生成得更完全（一个二醇对应两个醇分子），
且易于水解。
\textbf{6.} $15.09/150.9 = \qty{0.1000}{mol}$.
\textbf{7.} $1.20 \times 0.1000 \times 62.0 = \qty{7.44}{g}$.
\textbf{8.} \ce{BrCH2CH2CH2CHO + HOCH2CH2OH <=> C6H11BrO2 + H2O}
（2-（3-溴丙基）-1,3-二氧戊环）。
\textbf{9.} 它活化羰基，使第二阶段中
的失水得以进行；它被再生。
\textbf{10.} 甲苯与水一同蒸出，冷凝后落入
收集器；水下沉留存，甲苯溢流回去。
\textbf{11.} $Q$ 的分子中含 $[\ce{H2O}]$；除去水使
$Q < K$，反应继续进行。
\textbf{12.} 水的密度较大，且与甲苯不互溶：它聚集在
底部；只有上层液体才能到达回流臂。
\textbf{13.} 水不再聚集，即达到预期体积。
\textbf{14.} $0.1000 \times 194.9 = \qty{19.5}{g}$.
\textbf{15.} \ce{C6H11BrO2 + Mg -> C6H11BrMgO2}（于无水乙醚中）。
\textbf{16.} 该试剂加成到 \ce{(CH3)2CO} 上：生成\omterm{def:b1:alcohol-activation:alkoxide}{醇盐}
\ce{(CH3)2C(O^-)CH2CH2CH2-} 位于被\omterm{def:b1:acetals-protection:protecting-group}{保护}的链上，并带 \ce{MgBr+}。
\textbf{17.} 酸会过早除去\omterm{def:b1:acetals-protection:protecting-group}{保护基}，
并可能使叔醇\omterm{def:b1:alcohol-activation:dehydration}{脱水}。
\textbf{18.} $0.0700 \times 174.0 = \qty{12.2}{g}$.
\textbf{19.} \omterm{def:b1:acetals-protection:hemiacetal}{缩醛} + \ce{H2O} $\to$ 醛 + 乙烷-1,2-二醇
（酸催化）。
\textbf{20.} 其失水需要\omterm{def:b1:predominant-reaction:strong-acid}{强酸}和加热；温和而低温的
酸水溶液水解\omterm{def:b1:acetals-protection:hemiacetal}{缩醛}要快得多。
\textbf{21.} 5 号碳上的羟基加到醛基碳 1 上：一个六元环
（由五个碳和一个氧组成），与环氧相邻的碳上带两个甲基，
另一与之相邻的碳上带一个羟基。
\textbf{22.} $0.0700 \times 0.90 = \qty{0.0630}{mol}$, \qty{8.19}{g}.
\textbf{23.} $0.0630/0.1000 = \qty{63}{\percent}$（假定
\omterm{def:b1:acetals-protection:protecting-group}{保护}完全）。
\textbf{24.} $0.1000 \times 18.0 = \qty{1.80}{g}$.
\textbf{25.} $1.80/0.997 =$ \textbf{\qty{1.8}{mL} 的水}在捕集器中。
\end{solution}
